\documentclass[10pt, letterpaper]{article}

% Packages:
\usepackage[
    ignoreheadfoot, % set margins without considering header and footer
    top=2 cm, % seperation between body and page edge from the top
    bottom=2 cm, % seperation between body and page edge from the bottom
    left=1.5 cm, % seperation between body and page edge from the left
    right=1.5 cm, % seperation between body and page edge from the right
    footskip=1.0 cm, % seperation between body and footer
    % showframe % for debugging 
]{geometry} % for adjusting page geometry
\usepackage[explicit]{titlesec} % for customizing section titles
\usepackage{tabularx} % for making tables with fixed width columns
\usepackage{array} % tabularx requires this
\usepackage[dvipsnames]{xcolor} % for coloring text
\definecolor{primaryColor}{RGB}{53, 122, 178} % define primary color
\definecolor{secondaryColor}{RGB}{100, 100, 100} % define secondary color
\usepackage{enumitem} % for customizing lists
\usepackage{fontawesome5} % for using icons
\usepackage{amsmath} % for math
\usepackage[
    pdftitle={Baptiste Dubillaud's CV},
    pdfauthor={Baptiste Dubillaud},
    pdfcreator={LaTeX with RenderCV},
    colorlinks=true,
    urlcolor=primaryColor
]{hyperref} % for links, metadata and bookmarks
\usepackage[pscoord]{eso-pic} % for floating text on the page
\usepackage{calc} % for calculating lengths
\usepackage{bookmark} % for bookmarks
\usepackage{lastpage} % for getting the total number of pages
\usepackage{changepage} % for one column entries (adjustwidth environment)
\usepackage{paracol} % for two and three column entries
\usepackage{ifthen} % for conditional statements
\usepackage{needspace} % for avoiding page brake right after the section title
\usepackage{iftex} % check if engine is pdflatex, xetex or luatex

% Ensure that generate pdf is machine readable/ATS parsable:
\ifPDFTeX
    \input{glyphtounicode}
    \pdfgentounicode=1
    \usepackage[T1]{fontenc}
    \usepackage[utf8]{inputenc}
    \usepackage{lmodern}
\fi

\usepackage[default, type1]{sourcesanspro} 

% Some settings:
\AtBeginEnvironment{adjustwidth}{\partopsep0pt} % remove space before adjustwidth environment
\pagestyle{empty} % no header or footer
\setcounter{secnumdepth}{0} % no section numbering
\setlength{\parindent}{0pt} % no indentation
\setlength{\topskip}{0pt} % no top skip
\setlength{\columnsep}{0.15cm} % set column seperation
\makeatletter
\let\ps@customFooterStyle\ps@plain % Copy the plain style to customFooterStyle
\patchcmd{\ps@customFooterStyle}{\thepage}{
    \color{gray}\textit{\small Page \thepage{} of \pageref*{LastPage}}
}{}{} % replace number by desired string
\makeatother
\pagestyle{customFooterStyle}

\titleformat{\section}{
    % avoid page braking right after the section title
    \needspace{4\baselineskip}
    % make the font size of the section title large and color it with the primary color
    \Large\color{primaryColor}
}{
}{
}{
    % print bold title, give 0.15 cm space and draw a line of 0.8 pt thickness
    % from the end of the title to the end of the body
    \textbf{#1}\hspace{0.15cm}\titlerule[0.8pt]\hspace{-0.1cm}
}[] % section title formatting

\titlespacing{\section}{
    % left space:
    -1pt
}{
    % top space:
    0.5 cm
}{
    % bottom space:
    0.3 cm
} % section title spacing

% \renewcommand\labelitemi{$\vcenter{\hbox{\small$\bullet$}}$} % custom bullet points
\newenvironment{highlights}{
    \begin{itemize}[
        topsep=0.10 cm,
        parsep=0.10 cm,
        partopsep=0pt,
        itemsep=0pt,
        leftmargin=0.4 cm + 10pt
    ]
}{
    \end{itemize}
} % new environment for highlights

\newenvironment{highlightsforbulletentries}{
    \begin{itemize}[
        topsep=0.10 cm,
        parsep=0.10 cm,
        partopsep=0pt,
        itemsep=0pt,
        leftmargin=10pt
    ]
}{
    \end{itemize}
} % new environment for highlights for bullet entries


\newenvironment{onecolentry}{
    
} % new environment for one column entries

\newenvironment{twocolentry}[2][]{
    \onecolentry
    \def\secondColumn{#2}
    \setcolumnwidth{\fill, 6 cm}
    \begin{paracol}{2}
}{
    \switchcolumn \raggedleft \secondColumn
    \end{paracol}
    \endonecolentry
} % new environment for two column entries

\newenvironment{threecolentry}[3][]{
    \onecolentry
    \def\thirdColumn{#3}
    \setcolumnwidth{1 cm, \fill, 4.5 cm}
    \begin{paracol}{3}
    {\raggedright #2} \switchcolumn
}{
    \switchcolumn \raggedleft \thirdColumn
    \end{paracol}
    \endonecolentry
} % new environment for three column entries

\newenvironment{cliententry}[3]{
    % Mission client à l'intérieur d'un poste continu (freelance) : indentée sous
    % son expérience parente. L'en-tête aligne les dates à droite avec \hfill et
    % non avec twocolentry, car paracol ne se niche pas et prend ses largeurs sur
    % \textwidth, pas sur le \linewidth réduit par l'indentation.
    \begin{adjustwidth}{0.5 cm}{0 cm}
    {\color{primaryColor}\textbf{#1}}\hfill{\color{secondaryColor}#3}\par
    \ifthenelse{\equal{#2}{}}{}{\vspace{0.05 cm}{\color{secondaryColor}\textit{#2}}\par}
}{
    \end{adjustwidth}
} % new environment for a client mission inside an experience

\newenvironment{header}{
    \setlength{\topsep}{0pt}\par\kern\topsep\centering\color{primaryColor}\linespread{1.5}
}{
    \par\kern\topsep
} % new environment for the header

\newcommand{\placelastupdatedtext}{% \placetextbox{<horizontal pos>}{<vertical pos>}{<stuff>}
  \AddToShipoutPictureFG*{% Add <stuff> to current page foreground
    \put(
        \LenToUnit{\paperwidth-2 cm-0.2 cm+0.05cm},
        \LenToUnit{\paperheight-1.0 cm}
    ){\vtop{{\null}\makebox[0pt][c]{
        % \small\color{gray}\textit{Last updated in September 2024}\hspace{\widthof{Last updated in September 2024}}
    }}}%
  }%
}%

% save the original href command in a new command:
\let\hrefWithoutArrow\href

% new command for external links:
\renewcommand{\href}[2]{\hrefWithoutArrow{#1}{\ifthenelse{\equal{#2}{}}{ }{#2 }\raisebox{.15ex}{\footnotesize \faExternalLink*}}}


\begin{document}
    \newcommand{\AND}{\unskip
        \cleaders\copy\ANDbox\hskip\wd\ANDbox
        \ignorespaces
    }
    \newsavebox\ANDbox
    \sbox\ANDbox{}

    \placelastupdatedtext
    \begin{header}
        \fontsize{30 pt}{30 pt}
        \textbf{BAPTISTE \color{black} DUBILLAUD}

        \color{black}
        \vspace{-0.4 cm}
        \LARGE{Freelance Tech-Lead \& AI Software Engineer}

        \vspace{0.3 cm}

        \color{primaryColor}
        \normalsize
        \mbox{{\footnotesize\faMapMarker*}\hspace*{0.13cm}Pau, France}%
        % \kern 0.25 cm%
        % \AND%
        % \kern 0.25 cm%
        % \mbox{\hrefWithoutArrow{mailto:baptiste.dubillaud@gmail.com}{{\footnotesize\faEnvelope[regular]}\hspace*{0.13cm}baptiste.dubillaud@gmail.com}}%
        % \kern 0.25 cm%
        % \AND%
        % \kern 0.25 cm%
        % \mbox{\hrefWithoutArrow{tel:+33 6 01 15 67 84}{{\footnotesize\faPhone*}\hspace*{0.13cm}+33 6 01 15 67 84}}%
        \kern 0.25 cm%
        \AND%
        \kern 0.25 cm%
        \mbox{\hrefWithoutArrow{https://www.dubillaudb.fr/}{{\footnotesize\faGlobe}\hspace*{0.13cm}dubillaudb.fr}}%
        \kern 0.25 cm%
        \AND%
        \kern 0.25 cm%
        \mbox{\hrefWithoutArrow{https://www.linkedin.com/in/baptiste-dubillaud/?locale=en_EN}{{\footnotesize\faLinkedinIn}\hspace*{0.13cm}baptiste-dubillaud}}%
    \end{header}

    \vspace{0.4 cm}

    \begin{onecolentry}
        \begin{center}
            Software engineer with six years of experience across energy, aerospace and the public sector, now working as a \textbf{Freelance Tech-Lead \& AI Software Engineer}. Since 2023 I have designed and shipped \textbf{generative-AI} products (AI assistants, knowledge-extraction platforms, shared LLM APIs) from architecture to production, leading the teams that build them. Available for new \textbf{engagements} in Pau, Bordeaux, Paris or remotely.
        \end{center}
    \end{onecolentry}

    \section{Experience}

        \begin{twocolentry}{\color{secondaryColor}April 2026 – Present}
                \textbf{Freelance Tech-Lead \& AI Software Engineer}

            \textit{Pau/Bordeaux/Paris, France}
        \end{twocolentry}

        \vspace{0.1 cm}

        As a Freelance Tech-Lead \& AI Software Engineer, I work with clients to design and build AI-based software solutions, with a focus on code quality, best practices and client satisfaction.

        I handle every aspect of software development, from gathering requirements to final delivery, by way of architecture design, development, testing and deployment. I work closely with clients to understand their needs and deliver solutions that fit their business goals.

        \vspace{0.25 cm}

        \begin{cliententry}{DTNum – DGFiP}{Tech Lead – GenAI Project}{April 2026 – Present}
            \begin{highlights}
                \item \textbf{Digital Transformation Directorate (DTNum)} of the \textbf{DGFiP} – the French Public Tax Administration.
                \item Design and development of an \textbf{API} exposing locally-hosted models (\textbf{LLM}, \textbf{VLM}, rerank, embedding…) to the DGFiP's applications:
                    \begin{highlights}
                        \item Raw access to the models.
                        \item Turnkey \textbf{Chat} and \textbf{RAG} solutions, plus various applications.
                    \end{highlights}
                \item Development of the authentication, authorization and usage management (token consumption) system.
                \item Python (FastAPI, OpenAI SDK, LangChain, LiteLLM, SQLAlchemy, Alembic), PostgreSQL, Qdrant, Redis, MinIO/S3, Docker, Kubernetes.
            \end{highlights}
        \end{cliententry}

        \vspace{0.2 cm}

        \begin{twocolentry}{\color{secondaryColor}April 2025 – April 2026}
            \textbf{ThinkDeep AI, Senior GEN AI software engineer}

            \textit{Bordeaux/Paris, France}
        \end{twocolentry}
        
        \begin{highlights}
                \item Currently working as a Senior GEN AI Software Engineer at ThinkDeep AI, a startup specialized in developing AI-powered software solutions.

                \item \textbf{DeepBrain} : platform extracting the knowledge from documents \& data. Designed and implemented a group feature, integrated with Azure Entra ID, enabling secure knowledge and assistants sharing. Improved the user experience with new UI components and better performance.

                \item \textbf{AI Assistant}: working for DGFiP (French tax administration) to build an AI platform for public agents. Development of administration features (users, prompts, alerts) and improved overall project code quality by refactoring and formalization of practices.

                \item Exhibitor at \textbf{NVIDIA GTC Europe} in Paris.
        \end{highlights}

        \vspace{0.2 cm}
        
        \begin{twocolentry}{\color{secondaryColor}March 2023 – March 2025}
            \textbf{TotalEnergies, Tech Lead \& Software Engineer}

            \textit{Esbjerg, Denmark}
        \end{twocolentry}
        \begin{highlights}
            \item Create, refine and maintain Proof of Concepts for the affiliate's Digital Laboratory.
            \item Structure the development environment to support growth by managing different Github repositories, Azure Landing Zone, Windows servers, and implementing SSL and SSO...
            \item Develop applications utilizing Generative AI, complex user interactions (chat, animations...), large data fetching and processing, SAP data integration, health and usage monitoring, architecture, and Azure.
            \item Achieved Top 10 finalist status twice in TotalEnergies' \textit{E\&P Best Innovators} awards.
            \item React (Next.js, Axios, Framer-Motion, MSAL), Python (FastAPI, Pandas, Numpy, Dash), Powershell, PostgreSQL, Azure (AD, Entra ID, Vector Database, Azure functions).
        \end{highlights}

        \vspace{0.2 cm}

        \begin{twocolentry}{\color{secondaryColor}October 2021 – February 2023}
            \textbf{Airbus Defense \& Space, Software Engineer}

            \textit{Toulouse, France}
        \end{twocolentry}
            \begin{highlights}
                \item On behalf of \textbf{Viveris}, served as a member of an Airbus team developing a software for managing ground/satellite communications (encryption, decryption, routing), correcting anomalies, and implementing new features.
                \item Developed a simulator for the control center and satellite to support integration and performance tests of the application.
                \item Created a monitoring application to centralize logs and track the status of all different instances of the software deployed on multiple servers (both active and redundant instances).
                \item  Java (Netty, Swing), C++ (Qt), Python (Squish), Redmine, Squash, Jenkins, Gitlab, Sonarqube.
            \end{highlights}
        

        \vspace{0.2 cm}

        \begin{twocolentry}{\color{secondaryColor}September 2020 – September 2021}
            \textbf{TotalEnergies, Software Engineer}

            \textit{Pau, France}
        \end{twocolentry}

        \begin{highlights}
            \item Member of the AVO geophysicists and geoscientists team, working on a new set of computation and visualization tools for the SISMAGE-CIG, the geophysics software of TotalEnergies.
            \subitem Java (Swing, AWT), C, C++, Fortran, Git, Gerrit, Sonar, Jira, Jenkins.
            \item Project manager for the development of a web application of geophysics tools. Managed 3 interns.
            \subitem Angular, ExpressJS, NodeJS, BitBucket.
        \end{highlights}
        

        \vspace{0.2 cm}

        \begin{twocolentry}{\color{secondaryColor}May 2021 – August 2021}
            \textbf{TotalEnergies, Software Developer Intern}

            \textit{Pau, France}
        \end{twocolentry}
        
        \begin{highlights}
            \item Design and development of a well water injection simulation software.
            \item Java (Java FX, Control FX), Bash, Python, Fortran, BitBucket.
        \end{highlights}

        \vspace{0.2 cm}

        \begin{twocolentry}{\color{secondaryColor}May 2020 – Aug. 2020}
            \textbf{TotalEnergies, Software Developer Intern}

            \textit{Pau, France}
        \end{twocolentry}

        \begin{highlights}
            \item Development of visualization tools from scratch on SISMAGE-CIG, the geophysics platform of TotalEnergies on behalf of the geophysicists/geoscientists AVO team.
            \item Java (Swing), Git, Gerrit.
        \end{highlights}

        \vspace{0.2 cm}

        \begin{twocolentry}{\color{secondaryColor}June 2018 – August 2019}
            \textbf{Conforama, Salesman, student job}

            \textit{Lescar, France}
        \end{twocolentry}

        \begin{highlights}
            \item Seller of household appliances (TV, computers, washing machine…). 15 hours a week.
        \end{highlights}

    \section{Education}
    
        \begin{threecolentry}{\textbf{MS}}{
            \color{secondaryColor}2016 – 2021
        }
            \textbf{CY-Tech (EISTI), Computer Science}
            \begin{highlights}
                \item \textbf{Coursework:} Software Architecture, Algorithms, Machine Learning, Databases, Web development...
                \item Specialization in Big Data frameworks trough Big Data Analytics classes.
                \item Spark, Pandas, Seaborn, D3.js.
            \end{highlights}
        \end{threecolentry}

        \begin{threecolentry}{\textbf{-}}{
            \color{secondaryColor}2020 – 2021
        }
            \textbf{University of La Coroña, HPC Course}
            \begin{highlights}
                \item \textbf{Coursework:} HPC Architecture, GPU and CPU architecture, Compilers, Multi-threading algorithms, Image processing.
                \item One year course in High Performance Computing (HPC).
                \item Final project: Development and optimization of an algorithm to filter data cubes from seismic imaging (Pangea II HPC at TotalEnergies).
                \item Multi-threading (OpenMP, MPI, Cuda, OpenCL), Image processing (OpenCV), Fortran, C, C++.
            \end{highlights}
        \end{threecolentry}

    % \section{Projects}
    
    %     \begin{twocolentry}{
    %         \href{https://github.com/sinaatalay/rendercv}{github.com/name/repo}
    %     }
    %         \textbf{Multi-User Drawing Tool}
    %         \begin{highlights}
    %             \item Developed an electronic classroom where multiple users can simultaneously view and draw on a "chalkboard" with each person's edits synchronized
    %             \item Tools Used: C++, MFC
    %         \end{highlights}
    %     \end{twocolentry}


    %     \vspace{0.2 cm}

    %     \begin{twocolentry}{
    %         \href{https://github.com/sinaatalay/rendercv}{github.com/name/repo}
    %     }
    %         \textbf{Synchronized Desktop Calendar}
    %         \begin{highlights}
    %             \item Developed a desktop calendar with globally shared and synchronized calendars, allowing users to schedule meetings with other users
    %             \item Tools Used: C\#, .NET, SQL, XML
    %         \end{highlights}
    %     \end{twocolentry}


    %     \vspace{0.2 cm}

    %     \begin{twocolentry}{
    %         2002
    %     }
    %         \textbf{Custom Operating System}
    %         \begin{highlights}
    %             \item Built a UNIX-style OS with a scheduler, file system, text editor, and calculator
    %             \item Tools Used: C
    %         \end{highlights}
    %     \end{twocolentry}
    
    \section{Technologies}
        
        \begin{onecolentry}
            \textbf{Generative AI:} 
            \subitem OpenAI SDK, LangChain, LiteLLM, self-hosted models (LLM, VLM, embedding, rerank),
            \subitem RAG, vector databases: Qdrant, Vespa
        \end{onecolentry}

        \vspace{0.2 cm}

        \begin{onecolentry}
            \textbf{Software:} 
            \subitem Python: FastAPI, SQLAlchemy, Alembic, Pandas, NumPy,
            \subitem Java: AWT/SWING/JavaFX, Netty, Spark,
            \subitem C: OpenMP, MPI,
            \subitem C++: Qt, some basis in CUDA
        \end{onecolentry}

        \vspace{0.2 cm}

        \begin{onecolentry}
            \textbf{Web:} 
            \subitem Javascript: React (Next.js, Framer-Motion, MSAL), Vue.js, Angular, Node.js,
            \subitem Visualization: Recharts, D3.js, Dash, Streamlit
        \end{onecolentry}

        \vspace{0.2 cm}

        \begin{onecolentry}
            \textbf{Data:} PostgreSQL, Qdrant, Vespa, Redis, MinIO/S3, SAP, PI-Vision
        \end{onecolentry}

        \vspace{0.2 cm}

        \begin{onecolentry}
            \textbf{Deployment:} 
            \subitem Docker, Kubernetes, Poetry,
            \subitem Azure: Azure AD, Entra ID, Landing Zone management, Azure functions, Github actions,
            \subitem Windows Server: IIS and windows servers apps
        \end{onecolentry}

        \vspace{0.2 cm}

        \begin{onecolentry}
            \textbf{VCS \& tools:} Git, Github, Gitlab, Gerrit, Jenkins, Sonarqube, Jira, Redmine, VS Code, IntelliJ IDEA, DataGrip
        \end{onecolentry}


    \section{References}

    \begin{center}
    References available \textbf{upon request}.
    \end{center}

\end{document}